\clearpage
\begin{proofbox}
\textbf{\textcolor{CProof}{证明一：初中相似法——先用几何，再补抛物线定义}}

\textbf{先约定正长度。}对非竖直焦点弦，本证明把右端记为 $A$、左端记为 $B$。
令 $r=AF$、$s=BF$，并令 $u=x_A-p/2>0$、$v=p/2-x_B>0$。
它们正是上一节图中的两条水平边。

\par\medskip
\textbf{第1步：只用两个直角三角形相似。}
从两端向水平线作垂线，垂足为 $H,K$。
因 $A,F,B$ 共线，两个直角三角形在 $F$ 处锐角相等，故
\[
 \boxed{\frac{FH}{FA}=\frac{FK}{FB}},\qquad \frac ur=\frac vs.
\]
负斜率时把图关于水平线翻折，水平距离与相似关系都不变。

\par\medskip
\textbf{第2步：用到焦点等于到准线，把水平边换成斜边。}
准线 $x=-p/2$，所以
\[
 r=x_A+\frac p2=p+u,\quad
 s=x_B+\frac p2=p-v;
 \qquad u=r-p,\quad v=p-s.
\]

\par\medskip
\textbf{第3步：代入相似比，倒数自然出现。}
\[
\begin{aligned}
 \frac{r-p}{r}&=\frac{p-s}{s},\\
 1-\frac pr&=\frac ps-1,\\
 p\left(\frac1r+\frac1s\right)&=2,\\
 \boxed{\frac1r+\frac1s}&=\boxed{\frac2p}.
\end{aligned}
\]
\textbf{为何要除以斜边？}相似三角形给的就是这个比，不是为了凑倒数。

\par\medskip
\textbf{补齐竖直弦：}此时 $x_A=x_B=p/2$，由定义 $r=s=p$，所以倒数和也为 $2/p$。
\end{proofbox}

\clearpage
\begin{proofbox}
\textbf{\textcolor{CProof}{证明二：最快余弦法——把相似比换成三角函数}}

本证明改按上下标号：$A$ 为上端、$B$ 为下端；交换标号不影响倒数和。
令 $r=AF>0$、$s=BF>0$，倾斜角 $0<\theta<\pi$。
沿两条反向射线走到端点，横坐标为
\[
 x_A=\frac p2+r\cos\theta,\qquad
 x_B=\frac p2-s\cos\theta.
\]
抛物线定义给出 $AF=x_A+p/2$、$BF=x_B+p/2$，故
\[
 r=p+r\cos\theta,\qquad s=p-s\cos\theta.
\]
移项后分别除以 $pr$、$ps$：
\[
 \frac1r=\frac{1-\cos\theta}{p},\qquad
 \frac1s=\frac{1+\cos\theta}{p}.
\]
相加，余弦项相消：
\[
 \boxed{\frac1r+\frac1s}
 =\frac{1-\cos\theta+1+\cos\theta}{p}
 =\boxed{\frac2p}.
\]
这与相似法是同一原理：斜线两侧有相同的倾斜程度，水平分量符号相反。
若已经记住焦半径角度公式，考场上可以直接从两个倒数相加开始。
\end{proofbox}

\clearpage
\begin{proofbox}
\textbf{\textcolor{CProof}{证明三：坐标韦达法——只求两段长度的和与积}}

设过焦点的直线为 $x=my+p/2$，$m\in\mathbb R$。
$m=0$ 包含竖直通径；水平轴线不在这一族中，但它只有一个交点，本来不属于有限两点弦。

\par\medskip
\textbf{第1步：联立抛物线，得到一元二次式。}
\[
 y^2=2p\left(my+\frac p2\right)
 \ \Longrightarrow\ y^2-2pmy-p^2=0.
\]
两根为 $y_1,y_2$，于是
\[
 y_1+y_2=2pm,\qquad y_1y_2=-p^2.
\]
\textbf{第2步：用定义写焦半径。}
由 $x_i=my_i+p/2$ 得
\[
 r=my_1+p,\qquad s=my_2+p.
\]
\textbf{第3步：要算 $(r+s)/(rs)$，就分别算和、积。}
\[
\begin{aligned}
 r+s&=m(y_1+y_2)+2p=2p(1+m^2),\\
 rs&=m^2y_1y_2+mp(y_1+y_2)+p^2\\
   &=-p^2m^2+2p^2m^2+p^2=p^2(1+m^2).
\end{aligned}
\]
\textbf{第4步：通分相除，参数自动消掉。}
\[
 \frac1r+\frac1s=\frac{r+s}{rs}
 =\frac{2p(1+m^2)}{p^2(1+m^2)}=\boxed{\frac2p}.
\]
倒数和对两端交换不敏感，正好由韦达给出的和、积来求，无需解出交点根式。
\end{proofbox}

\clearpage
\begin{proofbox}
\textbf{\textcolor{CProof}{证明四：参数法——两端参数互为负倒数}}

取抛物线上参数点 $P(t)=(pt^2/2,pt)$，代入即满足 $y^2=2px$。
设不同两点 $A=P(t_1)$、$B=P(t_2)$。
由两点坐标消元可得弦所在直线
\[
 x=\frac{t_1+t_2}{2}\,y-\frac p2t_1t_2.
\]
代入焦点 $F(p/2,0)$，得到充要条件
\[
 \boxed{t_1t_2=-1}.
\]
定义给出
\[
 r=\frac p2(1+t_1^2),\qquad s=\frac p2(1+t_2^2).
\]
于是 $t_2=-1/t_1$，$t_1\ne0$，
\[
\begin{aligned}
 \frac1r+\frac1s
 &=\frac2p\left(\frac1{1+t_1^2}+\frac1{1+1/t_1^2}\right)\\
 &=\frac2p\left(\frac1{1+t_1^2}+\frac{t_1^2}{1+t_1^2}\right)
 =\boxed{\frac2p}.
\end{aligned}
\]
同一分母下的两项恰好补成 $1$，与前面的双色条形图含义一致。

\par\medskip
\textbf{继续推出长度界限与最短弦。}由倒数和式可得
\[
 s=\frac{pr}{2r-p},\quad r>p/2;
 \qquad 2rs=p(r+s).
\]
又 $(r+s)^2\ge4rs=2p(r+s)$，除以正数 $r+s$ 得 $r+s\ge2p$。
等号当且仅当 $r=s=p$，即通径。
\end{proofbox}
